\zlabel{triptych:concise:themes:start}
\section{The Propers: Themes and Movement}

The people who hope to meet Christ with deeds of justice ask God even for the will to do them. Among them are people scarcely able to ask for help. Others have already received considerable gifts and have other people's welfare in their hands. The Collect belongs to both. The man at the door in Jesus' story has one very definite task, to stay awake for his master. How much he may need to ask of God to do it becomes apparent in the garden where Jesus will find Peter asleep.

At the Entrance one voice from Psalm 25 (24) speaks in the presence of enemies who might laugh at the person who trusts God. Augustine hears Christ speaking for his Church. The Church has reason to remember what happened when she relied on herself. Some of the people who once encouraged her desires may now ridicule her for lifting her soul to God. Bellarmine attends to the other people in the psalm, all those who have waited for God. He understands how one person can venture to rely on a promise whose fulfillment other people have experienced. A person taking up the psalm for the first time has their company.\footnote{Augustine, \textit{Exposition on Psalm} 24, §§1--3. Robert Bellarmine, \textit{Commentary on the Book of Psalms}, Psalm 24, explanations 1--3.}

In Isaiah the whole people can barely account for themselves. Their own sins have carried them away as the wind takes a fallen leaf. Yet their first question to God sounds like an accusation. How did he let them stray so far? Jerome hears in this prayer the voices of different members of one assembly. Someone still intent on laying the blame elsewhere can speak alongside someone ready to confess his sin. By the end both have reason to call God Father. They belong to him as surely as clay in a potter's hands.

In Psalm 80 (79) the people ask God to come and look at his vine. It was he who brought it out of Egypt. Augustine notices the order in which the people speak about themselves at the end of the strophe. Given life, they will call upon God. He can ask how anyone makes the first small effort to call on God, and find in the singers' own words the one to whom they owe it. Each time the refrain comes round, the congregation asks for the sight of God's face. Augustine thinks of the face people could look on at Calvary. It was easy to recognize the man's suffering there and still fail to see whom they had before them.

\newpage
Paul thanks God for his gifts to the Corinthians. Soon he will have to ask members of this church why they have attached themselves to different teachers. The men who instructed them depended on the grace of Christ quite as much as the people who listened. A gifted speaker can help the entire church. Claiming him for one's own party is one way to deprive others of the good he can do them. Chrysostom hears Paul repeating the name of Christ throughout his thanksgiving. People acting as though they belonged to Paul or Apollos have been hearing the name of the one who was crucified for them all.

Before the Gospel the assembly asks for salvation in a verse of Psalm 85 (84). Augustine understands the petition as a desire to see the Savior himself. The Jesus whose words the disciples hear on the Mount of Olives speaks as a master might who wished to find his household well cared for on his return. One servant has been put at the door. As Jesus finishes the story he tells everyone to watch. The disciples who accompanied Jesus to the garden will shortly have the chance to stay awake beside someone they love.

At the altar the people have something to give. Grain can be grown and ground by people doing work that others depend on. Bread can also be given to someone without the means to buy it. The bishops' commentary on the offerings asks the worshippers to consider what has already been given to them that these gifts should be theirs to offer. In the Preface the priest gives thanks to the Father for the Son who shared our condition. As Cyril of Jerusalem explains the Creed, the one who stood silent while others judged him will himself be judge. The Church awaiting him in glory already knows what he endured for her.\footnote{Cyril of Jerusalem, \textit{Catechetical Lecture} XV, §§1--3. U.S. Roman Missal, Preface I of Advent.}

At Communion a second verse of Psalm 85 speaks of fruit. For Bellarmine the earth bears Christ when Mary gives birth to her Son. Augustine asks how the earth can yield fruit in a person who once delighted in doing wrong. After Communion the worshippers must still decide what to do with a loaf of bread, what to offer from their earnings, whom they wish to benefit. They ask to learn to love what will last. If the priest gives the solemn Advent blessing, one petition concerns the race they have yet to run. They can already practice the charity he asks for them.
\zlabel{triptych:concise:themes:end}
\clearpage
